% =============================================================================
% Annexe A, formulaire : toutes les formules effectivement implémentées.
% Ce fichier est inclus par rapport_scr.tex mais reste lisible seul.
% =============================================================================

\section*{Annexe A, formulaire de la formule standard}
\addcontentsline{toc}{section}{Annexe A, formulaire de la formule standard}

\small
\setlength{\parskip}{4pt}
\setlength{\parindent}{0pt}
Cette annexe recense les formules effectivement implémentées, avec leur référence dans le
règlement délégué (UE) 2015/35. Sauf mention contraire, le SCR d'un sous-module est la perte
de fonds propres de base provoquée par le choc :
\[
\mathrm{SCR}_i \;=\; \max\bigl(0,\; \mathrm{FPB}_0 - \mathrm{FPB}_i\bigr),
\qquad
\mathrm{FPB} \;=\; A - \mathrm{BE} - \mathrm{MR} - \text{autres passifs},
\]
où $A$ est la valeur de marché de l'actif. Les provisions techniques sont intégralement
reprojetées sous chaque choc, ce qui rend le calcul net des actions du management.

% -----------------------------------------------------------------------------
\subsection*{A.1 Provisions techniques et courbe des taux}

\textbf{Best estimate} (art. 28) : $\mathrm{BE} = \sum_{t} \mathrm{FT}_t \, v(t)$, avec
$v(t) = (1+s_t)^{-t}$ le facteur d'actualisation de la courbe sans risque.

\textbf{Extrapolation Smith-Wilson} (régime actuel) : les prix zéro-coupon s'écrivent
$P(t) = e^{-\omega t} + \sum_j \zeta_j\, W(t,u_j)$ avec $\omega = \ln(1+\mathrm{UFR})$ et
\[
W(t,u) = e^{-\omega(t+u)}\Bigl[\alpha\min(t,u) - \tfrac12 e^{-\alpha\max(t,u)}
\bigl(e^{\alpha\min(t,u)} - e^{-\alpha\min(t,u)}\bigr)\Bigr].
\]
Le paramètre $\alpha$ est le plus petit réel supérieur à $0{,}05$ tel que l'écart entre le
forward instantané au point de convergence et $\omega$ n'excède pas un point de base.

\textbf{Extrapolation alternative} (régime 2027) : au-delà du premier point de lissage
$\mathrm{FSP}$, le forward moyen sur $h$ années vaut
\[
f(h) = \mathrm{LLFR} + \bigl(\mathrm{UFR}^{c} - \mathrm{LLFR}\bigr)\,
\frac{1-e^{-a h}}{a h},
\]
où $\mathrm{LLFR}$ est la moyenne pondérée des forwards liquides et $\mathrm{UFR}^{c}$ l'UFR
en composition continue. Sous choc de taux, l'UFR est décalée de $\pm 15$ points de base.

\textbf{Marge de risque} (art. 37 à 39) :
\[
\mathrm{MR} = \mathrm{CoC} \sum_{t\ge 0} \frac{\mathrm{SCR}^{\text{nc}}(t)}{(1+s_{t+1})^{t+1}},
\qquad
\mathrm{SCR}^{\text{nc}}(t) \approx \mathrm{SCR}^{\text{nc}}(0)\,\frac{\mathrm{BE}(t)}{\mathrm{BE}(0)} ,
\]
où $\mathrm{SCR}^{\text{nc}}$ exclut le risque de marché. Le régime 2027 remplace
$\mathrm{CoC}=6\,\%$ par $4{,}75\,\%$ et pondère chaque année par
$\max(\lambda^{t}, \underline{\lambda})$.

% -----------------------------------------------------------------------------
\subsection*{A.2 Risque de marché (art. 164 à 188)}

\textbf{Taux} (art. 166 et 167). Régime actuel : $s_t^{\uparrow} = \max\{s_t(1+f^{\uparrow}_t),\,
s_t + 1\,\%\}$ et $s_t^{\downarrow} = s_t(1-f^{\downarrow}_t)$, sans choc sur les taux négatifs.
Régime 2027 : $s_t^{\uparrow} = s_t(1+b_t) + a_t$ et
$s_t^{\downarrow} = \max\{s_t(1-b_t) - a_t,\, \underline{s}_t\}$, le plancher $\underline{s}_t$
dépendant du terme.

\textbf{Action} (art. 168 à 172) : chocs de $39\,\%$ (type 1) et $49\,\%$ (type 2), majorés de
l'ajustement symétrique $\mathrm{SA}$ ; $22\,\%$ pour les participations stratégiques.
Agrégation :
$\mathrm{SCR}_{\text{act}} = \sqrt{\mathrm{SCR}_1^2 + 2\rho\,\mathrm{SCR}_1\mathrm{SCR}_2
+ \mathrm{SCR}_2^2}$ avec $\rho = 0{,}75$.

\textbf{Immobilier} (art. 174) : choc de $25\,\%$. \quad
\textbf{Change} (art. 188) : $\pm 25\,\%$ par devise, le pire scénario étant retenu.

\textbf{Spread} (art. 176 et 180) : pour une obligation de duration modifiée $d$ et d'échelon
de qualité de crédit $q$,
\[
\mathrm{stress}(q,d) = \min\bigl(a_{q,k} + b_{q,k}\,(d - \underline{d}_k),\; 1\bigr),
\qquad d \ge 1,
\]
$k$ désignant la bande de duration ($\le 5$, $]5,10]$, $]10,15]$, $]15,20]$, $>20$). Les
expositions souveraines de l'EEE libellées en monnaie locale sont exemptées.

\textbf{Concentration} (art. 184 à 187) : pour chaque groupe d'émetteurs $i$,
\[
\mathrm{XS}_i = \max\Bigl(0,\; \frac{E_i}{A^{xl}} - \mathrm{CT}_{q}\Bigr),
\qquad
\mathrm{Conc}_i = A^{xl}\,\mathrm{XS}_i\, g_{q},
\qquad
\mathrm{SCR}_{\text{conc}} = \sqrt{\textstyle\sum_i \mathrm{Conc}_i^2},
\]
$A^{xl}$ étant l'actif hors contrats en unités de compte.

\textbf{Agrégation} (art. 164) :
$\mathrm{SCR}_{\text{mar}} = \sqrt{\mathbf{v}^{\top} \mathbf{C}(A)\, \mathbf{v}}$, où le
paramètre $A$ vaut $0$ si le scénario de taux retenu est la hausse et $0{,}5$ s'il s'agit de la
baisse. Le régime 2027 ramène la corrélation taux-spread à $0{,}25$ en scénario baissier.

% -----------------------------------------------------------------------------
\subsection*{A.3 Défaut de la contrepartie (art. 189 à 202)}

Expositions de type 1 : $\mathrm{LGD}_j = 50\,\% \times (\text{recouvrables} + \text{créances})$
pour la réassurance. La variance de la perte agrégée vaut
\[
V = \sum_{j,k} \frac{\mathrm{PD}_j \mathrm{PD}_k}
{1{,}25(\mathrm{PD}_j+\mathrm{PD}_k) - \mathrm{PD}_j\mathrm{PD}_k}\,
\mathrm{TLGD}_j \mathrm{TLGD}_k
\;+\; \sum_j \frac{1{,}5\,\mathrm{PD}_j(1-\mathrm{PD}_j)}{2{,}5-\mathrm{PD}_j}\,
\sum_{i \in j}\mathrm{LGD}_i^2 ,
\]
puis
\[
\mathrm{SCR}_1 =
\begin{cases}
3\sqrt{V} & \text{si } \sqrt{V} \le 5\,\% \sum \mathrm{LGD},\\
5\sqrt{V} & \text{si } 5\,\% < \sqrt{V} \le 20\,\% \sum \mathrm{LGD},\\
\sum \mathrm{LGD} & \text{sinon.}
\end{cases}
\]
Type 2 : $\mathrm{SCR}_2 = 15\,\%\times$ créances non échues $+\; 90\,\%\times$ créances échues
depuis plus de trois mois. Agrégation :
$\mathrm{SCR}_{\text{déf}} = \sqrt{\mathrm{SCR}_1^2 + 1{,}5\,\mathrm{SCR}_1\mathrm{SCR}_2
+ \mathrm{SCR}_2^2}$.

% -----------------------------------------------------------------------------
\subsection*{A.4 Souscription vie et santé (art. 136 à 163)}

\begin{table}[H]\centering\footnotesize
\begin{tabularx}{\linewidth}{lXl}
\toprule
\textbf{Sous-module} & \textbf{Choc appliqué} & \textbf{Article}\\
\midrule
Mortalité & $q_x \rightarrow 1{,}15\,q_x$, permanent & 137\\
Longévité & $q_x \rightarrow 0{,}80\,q_x$, permanent & 138\\
Invalidité (santé) & taux de sortie d'incapacité réduits de $20\,\%$ & 153\\
Rachat & $\times 1{,}5$ ; $\times 0{,}5$ ; ou rachat massif de $40\,\%$ & 142\\
Frais & $+10\,\%$ de frais et $+1$ point d'inflation & 143\\
Révision & $+3\,\%$ (vie) ou $+4\,\%$ (santé) sur les arrérages & 141, 158\\
Catastrophe & $q_x \rightarrow q_x + 0{,}15$ point la première année & 144\\
\bottomrule
\end{tabularx}
\end{table}

Le choc de rachat massif n'est appliqué qu'aux contrats dont le rachat accroît les provisions,
contrat par contrat. L'agrégation utilise les matrices de l'annexe IV, de forme
$\sqrt{\mathbf{v}^{\top}\mathbf{C}\,\mathbf{v}}$. Le module santé combine ensuite SLT, NSLT et
catastrophe avec les corrélations $\rho_{\text{SLT,NSLT}} = 0{,}5$ et
$\rho_{\cdot,\text{CAT}} = 0{,}25$.

% -----------------------------------------------------------------------------
\subsection*{A.5 Souscription non-vie (art. 114 à 135)}

Pour chaque ligne d'activité $s$, avec $V_{p,s} = \max(P_s, P_{\text{last},s}) + \mathrm{FP}_s$
et $V_{r,s}$ le best estimate de sinistres :
\[
\sigma_s = \frac{\sqrt{(\sigma_{p,s}V_{p,s})^2 + \sigma_{p,s}V_{p,s}\sigma_{r,s}V_{r,s}
+ (\sigma_{r,s}V_{r,s})^2}}{V_{p,s}+V_{r,s}},
\qquad
\sigma = \frac{\sqrt{\sum_{s,t}\mathrm{Corr}_{s,t}\,\sigma_s V_s \sigma_t V_t}}{V},
\]
puis $\mathrm{SCR}_{\text{prim,rés}} = 3\,\sigma\,V$ avec $V = \sum_s V_s$.

\textbf{Paramètres propres à l'entreprise} (art. 218 et annexe XVII) :
$\sigma^{\star} = c\,\sigma^{\mathrm{USP}} + (1-c)\,\sigma^{\text{std}}$, le facteur de
crédibilité $c$ dépendant du nombre d'années d'historique ($c = 0{,}74$ pour dix ans).

\textbf{Catastrophe} : agrégation quadratique des périls naturels et des scénarios d'origine
humaine, puis application du traité en excédent de sinistres,
$L^{\text{net}} = \min(L, \text{priorité}) + \max(0,\, L - \text{priorité} - \text{portée})$.

% -----------------------------------------------------------------------------
\subsection*{A.6 Agrégation finale (art. 87, 103, 204 à 207, 248 à 253)}

\[
\mathrm{BSCR} = \sqrt{\sum_{i,j}\mathrm{Corr}_{i,j}\,\mathrm{SCR}_i\,\mathrm{SCR}_j}
\;+\; \mathrm{SCR}_{\text{incorporels}},
\]
\[
\mathrm{SCR}_{\text{op}} = \min\bigl(30\,\%\times\mathrm{BSCR},\;
\max(\mathrm{Op}_{\text{primes}}, \mathrm{Op}_{\text{prov}})\bigr)
+ 25\,\%\times \mathrm{Dép}_{\text{UC}},
\]
\[
\mathrm{SCR} = \mathrm{BSCR} + \mathrm{SCR}_{\text{op}}
+ \mathrm{Adj}_{\mathrm{TP}} + \mathrm{Adj}_{\mathrm{DT}},
\]
avec $\mathrm{Adj}_{\mathrm{TP}} = -\min\bigl(\mathrm{FDB},\;
\mathrm{BSCR}^{\text{brut}} - \mathrm{BSCR}^{\text{net}}\bigr)$ et
$\mathrm{Adj}_{\mathrm{DT}} = -\min\bigl(\tau\,(\mathrm{BSCR} + \mathrm{SCR}_{\text{op}}
+ \mathrm{Adj}_{\mathrm{TP}}),\; \text{capacité justifiable}\bigr)$, $\tau$ étant le taux
d'impôt. Enfin
\[
\mathrm{MCR} = \max\Bigl\{\min\bigl(\max(\mathrm{MCR}_{\text{lin}},\,25\,\%\,\mathrm{SCR}),\;
45\,\%\,\mathrm{SCR}\bigr),\; \mathrm{AMCR}\Bigr\}.
\]

% -----------------------------------------------------------------------------
\subsection*{A.7 Générateur de scénarios et modèle interne}

\textbf{Hull-White} : $\mathrm{d}r_t = (\theta(t) - a r_t)\mathrm{d}t + \sigma \mathrm{d}W_t$,
de solution exacte à pas annuel
$r_{t+1} = r_t e^{-a} + \alpha(t+1) - \alpha(t)e^{-a}
+ \sigma\sqrt{\tfrac{1-e^{-2a}}{2a}}\,\varepsilon$, avec
$\alpha(t) = f(0,t) + \tfrac{\sigma^2}{2a^2}(1-e^{-at})^2$, et prix zéro-coupon
$P(t,T) = A(t,T)e^{-B(t,T)r_t}$, $B(t,T) = (1-e^{-a(T-t)})/a$. Le contrôle market-consistent
est $\mathbb{E}[D_T] = P(0,T)$ et $\mathbb{E}[D_T S_T]/S_0 = 1$.

\textbf{Valeur temps des options} : $\mathrm{TVOG} = \mathrm{BE}^{\text{stoch}}
- \mathrm{BE}^{\text{central}}$ ; les prestations discrétionnaires futures s'obtiennent par
$\mathrm{FDB} = \mathrm{BE}^{\text{stoch}} - \mathrm{BE}^{\text{garanti}}$.

\textbf{Modèle interne} : $\mathrm{VaR}_{99{,}5\%}(L) - \mathbb{E}[L]$, la perte agrégée $L$
étant construite par copule, $L = \sum_k F_k^{-1}(U_k)$ avec $(U_1,\dots,U_n)$ gaussienne ou de
Student. La TVaR, $\mathbb{E}[L \mid L \ge \mathrm{VaR}]$, est calculée en parallèle pour
mesurer l'épaisseur de la queue.

\setlength{\parskip}{0pt}
\normalsize
